\chapter{Conclusiones y líneas futuras}
\label{cap:conclusiones}

\section{Conclusiones}

El proyecto ha alcanzado su objetivo general: se ha diseñado e implementado una aplicación web
que transforma vídeos de YouTube en material de estudio estructurado y consultable. Repasando los
objetivos específicos planteados en la introducción:

\begin{itemize}
  \item Se obtiene la transcripción de forma eficiente, priorizando los subtítulos de YouTube y
        recurriendo a Whisper solo cuando es necesario.
  \item El contenido se indexa mediante \textit{embeddings}, lo que habilita la búsqueda semántica
        por fragmentos.
  \item Se generan automáticamente capítulos y conceptos clave con marcas de tiempo coherentes.
  \item El asistente conversacional responde preguntas sobre el vídeo con memoria corta del
        diálogo y referencias al minuto exacto.
  \item Las clases se organizan en asignaturas y el sistema sugiere automáticamente la asignatura
        de cada vídeo nuevo.
  \item La interfaz integra reproductor, capítulos, conceptos y chat en una sola pantalla.
\end{itemize}

Más allá del cumplimiento de los requisitos, el trabajo me ha permitido integrar en un sistema
real un conjunto de tecnologías actuales —transcripción automática, \textit{embeddings},
búsqueda semántica y modelos de lenguaje— y enfrentarme a problemas propios de la ingeniería del
software, como el procesamiento asíncrono, la degradación controlada ante fallos de servicios
externos y la consistencia de los datos.

\subsection{Orientación social}

ComprendiA contribuye a un acceso más equitativo al conocimiento: permite navegar las clases por
conceptos en lugar de verlas completas, favorece el aprendizaje autónomo y ayuda a estudiantes
con ritmos o necesidades distintas. En este sentido, el proyecto se alinea con los Objetivos de
Desarrollo Sostenible relativos a la educación de calidad (ODS~4), la reducción de las
desigualdades (ODS~10) y la innovación (ODS~9).

\section{Líneas futuras}

El sistema es un prototipo funcional con margen claro de evolución. Las principales líneas de
trabajo futuro son:

\begin{itemize}
  \item \textbf{Integración con plataformas educativas (LMS).} Conectar ComprendiA con Moodle,
        Blackboard u otros sistemas para procesar directamente las clases del entorno
        universitario.
  \item \textbf{Asistente en tiempo real.} Llevar también el chat a WebSockets, con respuesta en
        \textit{streaming} (token a token), para una experiencia conversacional más inmediata. El
        progreso del análisis ya se transmite por WebSocket; faltaría extenderlo al asistente.
  \item \textbf{Gestión multiusuario.} Añadir autenticación y perfiles para soportar varios
        estudiantes y profesores.
  \item \textbf{Optimización del procesamiento.} Reducir la latencia de las llamadas externas
        (reutilizar la consulta de \texttt{yt-dlp}, cachear \textit{embeddings} de asignaturas,
        paralelizar fases) a partir de los tiempos medidos.
  \item \textbf{Mejora de la clasificación.} Umbral de similitud configurable, sugerencia también
        del profesor y refinamiento del aprendizaje del canal.
  \item \textbf{Reproductor enriquecido.} Marcar sobre la barra de progreso los capítulos y
        fragmentos analizados, sincronizando la posición del vídeo con la transcripción.
\end{itemize}

% TODO: si procede, mencionar el contexto del programa TALENT 25-26 y el recorrido del proyecto.
